\chapter{Requirements}
\ChapterLabel{Requirements}
\index{requirement}
\index{requirement!implementation version}
\index{priority}

\section{How to Use This Chapter}
State singular, necessary, unambiguous, feasible, and verifiable obligations.
Every requirement shall have a stable ID, one priority, the version in which it
was implemented, at least one mapped use case, and at least one mapped test case.
Enter \NoVersion\ in \textbf{In Version} when no implementation version has
been assigned. Preserve IDs when wording or version assignments change.

Select the requirement categories that make the specification easier to review
and govern. Keep each requirement in one primary register and use traceability
links instead of copying the same requirement into multiple categories.

\section{Requirement Identifier Naming Convention}
\label{Section:RequirementIdentifierConvention}
\index{requirement!identifier convention}
\index{requirement!family code}
\index{requirement!scope qualifier}

Requirement identifiers use the following grammar:
\begin{verbatim}
REQ-[<scope>-]<family>-<sequence>
\end{verbatim}
The optional scope qualifier identifies a project-defined subsystem, product,
content type, or other governed area. Omit it for project-wide requirements.
The sequence is a three-digit, zero-padded number that is unique within the
exact scope and family combination. For example, \texttt{REQ-F-004} is the
fourth project-wide functional requirement, while
\texttt{REQ-API-Q-003} is the third quality requirement in a project-defined
\texttt{API} scope.

\small
\begin{tabularx}{\linewidth}{L{2.0cm}L{3.5cm}Y}
\caption{Requirement Identifier Segments and Family Codes\label{Table:RequirementIdentifierConvention}}\\
\toprule
\textbf{Code} & \textbf{Meaning} & \textbf{Use}\\
\midrule
\endfirsthead
\multicolumn{3}{l}{\small\itshape Requirement Identifier Segments and Family Codes (continued)}\\
\toprule
\textbf{Code} & \textbf{Meaning} & \textbf{Use}\\
\midrule
\endhead
\bottomrule
\endlastfoot
\texttt{REQ} & Requirement & Identifies the artifact as a requirement. Every requirement ID begins with this segment.\\
\texttt{<scope>} & Project-defined scope & Optional uppercase qualifier for a governed subsystem, product, content type, or other area. Define every scope code before use; absence denotes project-wide scope.\\
\texttt{B} & Business & States the stakeholder outcome, mission, commercial purpose, policy goal, or business result the system shall enable.\\
\texttt{F} & Functional & States observable behavior, service, workflow, calculation, state change, or capability the system shall provide.\\
\texttt{Q} & Quality & States a measurable quality attribute or non-functional expectation, including accessibility, performance, reliability, resilience, usability, maintainability, or compatibility.\\
\texttt{I} & Interface & States a user, software, API, hardware, communication, protocol, or external-system boundary contract, including integration behavior at that boundary.\\
\texttt{S} & Security & States security, privacy, consent, authorization, auditing, secret-handling, or protected-data obligations maintained in a specialized security register.\\
\texttt{D} & Data & States data definition, authority, provenance, quality, integrity, ownership, lifecycle, retention, or source-of-truth obligations.\\
\texttt{C} & Constraint & States a binding implementation, platform, technology, format, storage, source-control, standards, compatibility, schedule, budget, or governance constraint.\\
\texttt{\#\#\#} & Sequence & Three-digit, zero-padded stable sequence within the exact scope and family; assigned monotonically and never reused.\\
\bottomrule
\end{tabularx}
\normalsize

Classify a requirement by its primary obligation and verification concern. A
requirement may support several concerns, but its identifier contains only one
family code. Represent cross-cutting relationships through traceability rather
than a compound family code.

An AI agent shall reuse an existing scope and family whenever their definitions
fit. If no defined family accurately represents the requirement, the agent
shall not silently invent or use a new code. It shall propose a unique uppercase
code, full family name, decision boundary, rationale, and representative
examples; check the complete requirements register and category catalog for
semantic overlap and identifier collision; and obtain authorized-human
approval. After approval, the agent shall add the family to
Table~\ref{Table:RequirementIdentifierConvention}, update the category catalog
and AI guidance, and create or update the applicable register before assigning
an identifier from it. The same approval and documentation rule applies to a
new scope qualifier. Extending the convention shall not rename, renumber, or
reuse an existing identifier.

\section{Functional and Quality Requirements}
\index{requirement!functional}
\index{requirement!quality}
Functional and quality requirements answer different questions and shall not be
used interchangeably.

\begin{dsnNote}[Central Distinction]
A \textbf{functional requirement} states what behavior, service, calculation,
state transition, or externally observable result the system shall provide. A
\textbf{quality requirement} states how well that behavior or the system as a
whole shall perform under stated conditions. Quality requirements require a
measurable threshold, workload, environment, or other evaluation basis.
\end{dsnNote}

For example, ``The system shall display the current request status'' is
functional. ``The system shall display the current request status within two
seconds under the reference workload'' is a quality requirement. Avoid vague
quality statements such as ``the system shall be fast, secure, reliable, or
easy to use.'' Replace them with measurable conditions and success criteria.

Quality requirements are also called quality-attribute or non-functional
requirements. The term \emph{quality requirement} is preferred here because
these obligations describe behavior that shall be designed, implemented, and
verified rather than an absence of function.

\section{Requirement Status Model}
\index{requirement!status}
\index{status!Open}
\index{status!In Progress}
\index{status!Done}
\index{status!Retired}
Priority and status are independent. Priority records the requirement's
importance or commitment; status records its progress through the delivery
lifecycle. A status change shall not erase or replace the assigned priority.

\begin{tabularx}{\linewidth}{L{3.0cm}YL{3.3cm}}
\caption{Requirement Status Definitions\label{Table:RequirementStatusDefinitions}}\\
\toprule
\textbf{Status} & \textbf{Meaning} & \textbf{In Version Rule}\\
\midrule
\endfirsthead
\toprule
\textbf{Status} & \textbf{Meaning} & \textbf{In Version Rule}\\
\midrule
\endhead
\ReqStatusOpen & The requirement is active, but implementation has not started. Deferred work remains Open until work begins or the requirement is retired. & Enter \NoVersion.\\
\ReqStatusInProgress & Implementation or integration work has started, but the requirement does not yet satisfy the definition of Done. & Enter \NoVersion.\\
\ReqStatusDone & The requirement is implemented; mapped tests pass; verification is complete; traceability is current; and the completed behavior is included in a controlled version. & Enter the first controlled version containing the completed requirement.\\
\ReqStatusRetired & The requirement is no longer active because it was superseded, made obsolete, rejected, or withdrawn. Retain its stable ID, original text, and traceability; record the rationale and any replacement in document history. & Retain the original implementation version if it was previously Done; otherwise enter \NoVersion.\\
\bottomrule
\end{tabularx}

Only \ReqStatusDone\ asserts completion. Do not remove a retired requirement or
reuse its ID. A requirement may move from \ReqStatusRetired\ only through an
approved change that restores its active status and records the decision in
document history.

\section{Requirement Category Catalog}
\index{requirement!category catalog}
The following catalog distinguishes broad requirement families from their
subtypes. In particular, performance, capacity, scalability, reliability,
availability, resilience, recoverability, usability, accessibility,
maintainability, supportability, testability, portability, compatibility,
interoperability, security, privacy, safety, sustainability, and observability
are quality attributes. A requirement concerning one of these attributes is a
quality requirement when it specifies how well the system shall operate under
stated conditions.

A project may organize a specialized quality attribute, such as security,
privacy, or safety, in a separate register when it has distinct owners, review
criteria, evidence, or assurance obligations. The separate register is a
management view; it does not change the requirement's classification as a
quality requirement. Likewise, interface, data, operational, and other families
may contain both functional obligations and associated quality obligations. Use
the primary purpose of each requirement to select its register and use links to
represent cross-cutting concerns.

The catalog is intentionally broad. A project may use a family as a top-level
register, create registers for selected subtypes, combine closely related
groups, or omit groups that do not apply. Category choice shall not weaken the
traceability or verification rules for individual requirements.

\section{Domain Profile Impact on Requirements}
\label{Section:DomainProfileRequirements}
\index{requirement!domain profile}
Every Applicable domain profile shall be evaluated when developing or changing
requirements. Each applicable profile concern shall result in one or more
project requirements, a reference to an existing requirement, or a documented
Not Applicable decision with rationale approved by an authorized human. Record
the profile applicability decision in the Development Plan and retain
project-specific obligations only in the applicable requirement registers.
The profile itself is a tailoring guide, not a duplicate requirements
specification. See Section~\ref{Section:DomainProfiles} and the applicable
profile appendix.

\small
\begin{tabularx}{\linewidth}{L{3.0cm}L{4.0cm}Y}
\caption{Requirement Families, Subtypes, and Quality Attributes\label{Table:RequirementCategoryCatalog}}\\
\toprule
\textbf{Requirement Family} & \textbf{Subtype or Quality Attribute} & \textbf{Typical Scope}\\
\midrule
\endfirsthead
\multicolumn{3}{l}{\small\itshape Requirement Families, Subtypes, and Quality Attributes (continued)}\\
\toprule
\textbf{Requirement Family} & \textbf{Subtype or Quality Attribute} & \textbf{Typical Scope}\\
\midrule
\endhead
\midrule
\multicolumn{3}{r}{\small\itshape Continued on next page}\\
\endfoot
\bottomrule
\endlastfoot
Business & Outcomes and policy & Organizational outcomes, value, policy goals, measurable benefits, cost targets, and mission objectives.\\
Stakeholder and user & Stakeholder need & Capabilities, outcomes, and constraints needed by a stakeholder group or external party.\\
Stakeholder and user & User need & User goals and externally visible outcomes expressed independently of a particular implementation.\\
Functional & System behavior & Services, calculations, decisions, state changes, workflows, and observable system responses.\\
Quality & General quality & Measurable characteristics that state how well a function or the system shall operate under stated conditions.\\
Quality & Performance & Response time, latency, throughput, processing time, and resource consumption.\\
Quality & Capacity & Supported data volume, transaction volume, concurrent users, devices, or retained records.\\
Quality & Scalability & Behavior as workload, data, users, or deployed instances increase or decrease.\\
Quality & Reliability & Probability or frequency of failure-free operation under stated conditions.\\
Quality & Availability & Required service availability, planned downtime, and permitted outage duration.\\
Quality & Resilience and fault tolerance & Continued or degraded operation when components, dependencies, or communication paths fail.\\
Quality & Recoverability and continuity & Recovery time, recovery point, restoration, backup, failover, and disaster-continuity objectives.\\
Quality & Usability & Task effectiveness, efficiency, learnability, error prevention, and user satisfaction.\\
Quality & Accessibility & Conformance and measurable support for users with visual, auditory, motor, speech, or cognitive disabilities.\\
Quality & Maintainability & Time, effort, modularity, analyzability, and change characteristics required for system modification.\\
Quality & Supportability & Diagnostic access, service procedures, field replacement, support hours, and support-team capabilities.\\
Quality & Testability & Controllability, observability, repeatability, test interfaces, and evidence needed to verify behavior.\\
Quality & Portability & Ability to install, move, or operate the system across specified environments or platforms.\\
Quality & Compatibility & Coexistence and correct operation with specified versions, devices, browsers, operating systems, or products.\\
Quality & Interoperability & Coordinated behavior, protocols, semantics, orchestration, and exchanges with external systems.\\
Quality & Security & Authentication, authorization, confidentiality, integrity, cryptography, auditing, and threat resistance.\\
Quality & Privacy & Collection, consent, purpose, minimization, access, disclosure, retention, deletion, and rights for personal data.\\
Quality & Safety & Hazard prevention, risk reduction, safe state, interlocks, alarms, and unacceptable-loss controls.\\
Quality & Sustainability & Energy use, resource efficiency, carbon reporting, hardware life, repair, reuse, and disposal objectives.\\
Quality & Observability & Logs, metrics, traces, health signals, diagnostic context, and measurable diagnostic coverage.\\
Interface & User interface & Human-system interaction, presentation, navigation, input, feedback, and interaction contracts.\\
Interface & Software and API & Operations, messages, schemas, errors, protocols, sequencing, and software boundary contracts.\\
Interface & Hardware and communication & Connectors, signals, electrical characteristics, device protocols, networks, and physical boundaries.\\
Data and information & Data definition and quality & Data definitions, schemas, provenance, accuracy, completeness, integrity, ownership, and permitted transformations.\\
Data and information & Lifecycle and governance & Creation, access, retention, archival, legal hold, disposition, deletion, and proof of destruction.\\
Operational & Operation and administration & Startup, shutdown, scheduling, configuration, provisioning, monitoring, backup, recovery, and incident response.\\
Operational & Deployment and platform & Approved platforms, topology, packaging, installation, configuration, upgrade, and runtime environment.\\
Operational & Infrastructure and environment & Network, compute, storage, facility, power, temperature, physical location, and environmental limits.\\
Transition & Migration and rollout & Data conversion, coexistence, rollout, cutover, rollback, legacy integration, and decommissioning.\\
Compliance and legal & Regulatory and standards & Obligations imposed by regulations, standards, certifications, policies, or assurance frameworks.\\
Compliance and legal & Legal and contractual & Licenses, warranties, intellectual-property terms, service agreements, procurement terms, and contractual duties.\\
Compliance and legal & Reporting and audit & Reports, evidence, audit trails, review access, content, timing, and record accountability.\\
Product support & Documentation & Required user, operator, administrator, developer, API, maintenance, and compliance documentation.\\
Product support & Training & Required curricula, exercises, qualifications, delivery modes, and demonstrated learner outcomes.\\
Commercial & Licensing and entitlement & Product editions, entitlement, metering, third-party components, subscriptions, and licensing restrictions.\\
AI and machine learning & Model and data controls & Model performance, data provenance, bias evaluation, explainability, human oversight, drift, and retraining controls.\\
Physical and hardware & Physical characteristics & Dimensions, mass, materials, connectors, tolerances, environmental ratings, and hardware behavior.\\
Constraint & Solution constraint & Mandated technologies, designs, budgets, schedules, standards, dependencies, or solution boundaries.\\
\end{tabularx}
\normalsize

\section{Creating a Project-Specific Category}
Create an additional category when it gives requirements a distinct meaning,
owner, review method, or assurance process. Define the category's scope and a
unique family code using
Section~\ref{Section:RequirementIdentifierConvention}; duplicate one of the
category-specific registers below; keep the standard columns; add the category to
\hyperref[Table:RequirementCategoryCatalog]{the category catalog}; and update
use cases, tests, traceability, verification, validation, document history, and
the index as applicable. Document whether the new category is top-level or a
subtype of an existing category.

Do not create a category solely to avoid resolving an ambiguous requirement.
If one statement contains multiple independently verifiable obligations, split
it into singular requirements and assign each one a primary category.

\section{Review Rules}
Reject a requirement row if its priority, status, use-case mapping, or test-case
mapping is empty. The \textbf{Status} field shall contain Open, In Progress,
Done, or Retired. The \textbf{In Version} field shall follow
\hyperref[Table:RequirementStatusDefinitions]{the status definitions}. Confirm
that reciprocal mappings appear in the user-story, use-case, test,
traceability, and verification registers.

\begin{landscape}
\small
\section{Category-Specific Requirement Registers}
The seven registers below are the template's default working set. They are
selected for practical review and ownership, not as seven mutually exclusive
top-level families. For example, the Security Requirements Register is a
specialized quality-requirement register. Add, rename, combine, or remove
registers to match the approved project taxonomy. Every retained register shall
use the same traceability columns.

\subsection{Business Requirements}
\index{requirement!business}
\begin{tabularx}{\linewidth}{L{2.1cm}YL{2.0cm}L{2.1cm}L{1.5cm}L{2.6cm}L{2.6cm}}
\caption{Business Requirements Register\label{Table:BusinessRequirements}}\\
\toprule
\textbf{ID} & \textbf{Requirement} & \textbf{Priority} & \textbf{Status} & \textbf{In Version} & \textbf{Use Case(s)} & \textbf{Test Case(s)}\\
\midrule
\endfirsthead
\multicolumn{7}{l}{\small\itshape Business Requirements Register (continued)}\\
\toprule
\textbf{ID} & \textbf{Requirement} & \textbf{Priority} & \textbf{Status} & \textbf{In Version} & \textbf{Use Case(s)} & \textbf{Test Case(s)}\\
\midrule
\endhead
\midrule
\multicolumn{7}{r}{\small\itshape Continued on next page}\\
\endfoot
\bottomrule
\endlastfoot
\RequirementID{REQ-B-001} & The system shall enable an authorized analyst to prepare approved request-history data for analysis within five minutes for the reference data set. & \gls{must} & \ReqStatusOpen & \NoVersion & \UseCaseRef{UC-003} & \TestCaseRef{TC-009}\\
\end{tabularx}

\clearpage
\subsection{Functional Requirements}
\begin{tabularx}{\linewidth}{L{2.1cm}YL{2.0cm}L{2.1cm}L{1.5cm}L{2.6cm}L{2.6cm}}
\caption{Functional Requirements Register\label{Table:FunctionalRequirements}}\\
\toprule
\textbf{ID} & \textbf{Requirement} & \textbf{Priority} & \textbf{Status} & \textbf{In Version} & \textbf{Use Case(s)} & \textbf{Test Case(s)}\\
\midrule
\endfirsthead
\multicolumn{7}{l}{\small\itshape Functional Requirements Register (continued)}\\
\toprule
\textbf{ID} & \textbf{Requirement} & \textbf{Priority} & \textbf{Status} & \textbf{In Version} & \textbf{Use Case(s)} & \textbf{Test Case(s)}\\
\midrule
\endhead
\midrule
\multicolumn{7}{r}{\small\itshape Continued on next page}\\
\endfoot
\bottomrule
\endlastfoot
\RequirementID{REQ-F-001} & The system shall accept a valid service request and assign it a unique identifier. & \gls{must} & \ReqStatusDone & 1.0 & \UseCaseRef{UC-001} & \TestCaseRef{TC-001}\\
\RequirementID{REQ-F-002} & The system shall display the current status and most recent update for a submitted request. & \gls{should} & \ReqStatusDone & 1.1 & \UseCaseRef{UC-002} & \TestCaseRef{TC-002}\\
\RequirementID{REQ-F-003} & The system shall export request history in the approved interchange format. & \gls{would} & \ReqStatusOpen & \NoVersion & \UseCaseRef{UC-003} & \TestCaseRef{TC-004}\\
\end{tabularx}

\subsection{Quality Requirements}
\begin{tabularx}{\linewidth}{L{2.1cm}YL{2.0cm}L{2.1cm}L{1.5cm}L{2.6cm}L{2.6cm}}
\caption{Quality Requirements Register\label{Table:QualityRequirements}}\\
\toprule
\textbf{ID} & \textbf{Requirement} & \textbf{Priority} & \textbf{Status} & \textbf{In Version} & \textbf{Use Case(s)} & \textbf{Test Case(s)}\\
\midrule
\endfirsthead
\multicolumn{7}{l}{\small\itshape Quality Requirements Register (continued)}\\
\toprule
\textbf{ID} & \textbf{Requirement} & \textbf{Priority} & \textbf{Status} & \textbf{In Version} & \textbf{Use Case(s)} & \textbf{Test Case(s)}\\
\midrule
\endhead
\midrule
\multicolumn{7}{r}{\small\itshape Continued on next page}\\
\endfoot
\bottomrule
\endlastfoot
\RequirementID{REQ-Q-001} & The system shall display a request status within two seconds under the reference workload. & \gls{could} & \ReqStatusDone & 1.2 & \UseCaseRef{UC-002} & \TestCaseRef{TC-003}\\
\end{tabularx}

\clearpage
\subsection{Interface Requirements}
\index{requirement!interface}
\begin{tabularx}{\linewidth}{L{2.1cm}YL{2.0cm}L{2.1cm}L{1.5cm}L{2.6cm}L{2.6cm}}
\caption{Interface Requirements Register\label{Table:InterfaceRequirements}}\\
\toprule
\textbf{ID} & \textbf{Requirement} & \textbf{Priority} & \textbf{Status} & \textbf{In Version} & \textbf{Use Case(s)} & \textbf{Test Case(s)}\\
\midrule
\endfirsthead
\multicolumn{7}{l}{\small\itshape Interface Requirements Register (continued)}\\
\toprule
\textbf{ID} & \textbf{Requirement} & \textbf{Priority} & \textbf{Status} & \textbf{In Version} & \textbf{Use Case(s)} & \textbf{Test Case(s)}\\
\midrule
\endhead
\midrule
\multicolumn{7}{r}{\small\itshape Continued on next page}\\
\endfoot
\bottomrule
\endlastfoot
\RequirementID{REQ-I-001} & The system shall exchange request-history exports using the approved, versioned interface schema. & \gls{should} & \ReqStatusInProgress & \NoVersion & \UseCaseRef{UC-003} & \TestCaseRef{TC-005}\\
\end{tabularx}

\subsection{Security Requirements}
\index{requirement!security}
\begin{tabularx}{\linewidth}{L{2.1cm}YL{2.0cm}L{2.1cm}L{1.5cm}L{2.6cm}L{2.6cm}}
\caption{Security Requirements Register\label{Table:SecurityRequirements}}\\
\toprule
\textbf{ID} & \textbf{Requirement} & \textbf{Priority} & \textbf{Status} & \textbf{In Version} & \textbf{Use Case(s)} & \textbf{Test Case(s)}\\
\midrule
\endfirsthead
\multicolumn{7}{l}{\small\itshape Security Requirements Register (continued)}\\
\toprule
\textbf{ID} & \textbf{Requirement} & \textbf{Priority} & \textbf{Status} & \textbf{In Version} & \textbf{Use Case(s)} & \textbf{Test Case(s)}\\
\midrule
\endhead
\midrule
\multicolumn{7}{r}{\small\itshape Continued on next page}\\
\endfoot
\bottomrule
\endlastfoot
\RequirementID{REQ-S-001} & The system shall permit request-history exports only to authenticated actors assigned the approved export role. & \gls{must} & \ReqStatusDone & 1.0 & \UseCaseRef{UC-003} & \TestCaseRef{TC-006}\\
\end{tabularx}

\clearpage
\subsection{Data Requirements}
\index{requirement!data}
\begin{tabularx}{\linewidth}{L{2.1cm}YL{2.0cm}L{2.1cm}L{1.5cm}L{2.6cm}L{2.6cm}}
\caption{Data Requirements Register\label{Table:DataRequirements}}\\
\toprule
\textbf{ID} & \textbf{Requirement} & \textbf{Priority} & \textbf{Status} & \textbf{In Version} & \textbf{Use Case(s)} & \textbf{Test Case(s)}\\
\midrule
\endfirsthead
\multicolumn{7}{l}{\small\itshape Data Requirements Register (continued)}\\
\toprule
\textbf{ID} & \textbf{Requirement} & \textbf{Priority} & \textbf{Status} & \textbf{In Version} & \textbf{Use Case(s)} & \textbf{Test Case(s)}\\
\midrule
\endhead
\midrule
\multicolumn{7}{r}{\small\itshape Continued on next page}\\
\endfoot
\bottomrule
\endlastfoot
\RequirementID{REQ-D-001} & The system shall retain each request status transition with its timestamp, actor identifier, previous status, and new status. & \gls{must} & \ReqStatusDone & 1.1 & \UseCaseRef{UC-002}, \UseCaseRef{UC-003} & \TestCaseRef{TC-007}\\
\end{tabularx}

\subsection{Constraint Requirements}
\index{requirement!constraint}
\begin{tabularx}{\linewidth}{L{2.1cm}YL{2.0cm}L{2.1cm}L{1.5cm}L{2.6cm}L{2.6cm}}
\caption{Constraint Requirements Register\label{Table:ConstraintRequirements}}\\
\toprule
\textbf{ID} & \textbf{Requirement} & \textbf{Priority} & \textbf{Status} & \textbf{In Version} & \textbf{Use Case(s)} & \textbf{Test Case(s)}\\
\midrule
\endfirsthead
\multicolumn{7}{l}{\small\itshape Constraint Requirements Register (continued)}\\
\toprule
\textbf{ID} & \textbf{Requirement} & \textbf{Priority} & \textbf{Status} & \textbf{In Version} & \textbf{Use Case(s)} & \textbf{Test Case(s)}\\
\midrule
\endhead
\midrule
\multicolumn{7}{r}{\small\itshape Continued on next page}\\
\endfoot
\bottomrule
\endlastfoot
\RequirementID{REQ-C-001} & The system shall support the deployment platform and browser versions identified in the approved deployment baseline. & \gls{must} & \ReqStatusDone & 1.0 & \UseCaseRef{UC-001}, \UseCaseRef{UC-002}, \UseCaseRef{UC-003} & \TestCaseRef{TC-008}\\
\end{tabularx}
\end{landscape}
